\documentclass[10pt,a4paper]{amsart}
\usepackage[margin=19mm]{geometry}
\usepackage{amsmath,amssymb,mathtools}
\usepackage[scr=rsfs,cal=euler]{mathalfa}
\usepackage{xfrac}
\usepackage[unicode,hidelinks,bookmarks=false]{hyperref}
\newcommand{\ie}{i.e.\ }
\newcommand{\cf}{cf.\ }
\newcommand{\resp}{respectively\ }
\newcommand{\Subsection}{\subsection}
\providecommand{\nobreakdash}{\nobreak\hbox{-}}
\setlength{\emergencystretch}{2em}
\multlinegap0pt
\usepackage{mathrsfs,enumitem}
\numberwithin{equation}{section}
\newtheorem{theorem}{Theorem}
\newtheorem{proposition}[theorem]{Proposition}
\newtheorem{lemma}[theorem]{Lemma}
\newtheorem{corollary}[theorem]{Corollary}
\newtheorem{remark}[theorem]{Remark}
\newtheorem{definition}[theorem]{Definition}
\newtheorem{conjecture}[theorem]{Conjecture}
\newtheorem{problem}[theorem]{Problem}
\newtheorem{example}[theorem]{Example}
\newcommand{\bC}{\mathbb C}
\newcommand{\bR}{\mathbb R}
\newcommand{\D}{\mathcal D}
\newcommand{\Res}{\operatorname{Res}}
\newcommand{\disc}{\operatorname{disc}}
\newcommand{\Sym}{\operatorname{Sym}}
\newcommand{\barx}{\overline{x}}
\begin{document}
\title[Discriminants of derivatives and symmetric difference polyno]{Discriminants of derivatives and symmetric difference polynomials}
\author{Boris Shapiro}
\date{}
\maketitle
\noindent\emph{Source and licence.} Discriminants of derivatives and symmetric difference polynomials. Version arXiv:2605.25743v2 (CC BY 4.0). This material is distributed under the Creative Commons Attribution 4.0 International licence, http://creativecommons.org/licenses/by/4.0/, as recorded in the arXiv OAI licence metadata for this exact version and shown on the abstract page https://arxiv.org/abs/2605.25743v2. Full rights notice: licenses/arxiv-2605.25743-CC-BY-4.0.txt.\par\smallskip
\noindent\textit{Editorial scope.} Counted unit: the original \texttt{theorem} environment \texttt{thm:terminal-normalized} at source lines 613--651 together with its complete proof at lines 653--681; the delivered document keeps source lines 604--682 so that every referenced label and every citation resolves inside it. Native editable source is the paper's own concatenated TeX; the AMS wrapper is editorial formatting only. No publisher class, upstream script or unrelated artwork is redistributed.
\section{Universal terminal polynomials}

For $m\ge j$ write
\[
(m)_{\downarrow j}=m(m-1)\cdots(m-j+1)
\]
for the falling factorial.  The preceding examples are instances of the
following normalized terminal principle.


\begin{theorem}[subset-average and Appell structure]\label{thm:terminal-normalized}
Fix $r\ge2$ and $n\ge r$.  After translating the roots so that
$\sum_i x_i=0$, put
\[
J_{n,r}(t)=\frac{r!}{n!}P^{(n-r)}(t).
\]
Then one has the subset-average representation
\begin{equation}\label{eq:subset-average}
J_{n,r}(t)=\binom nr^{-1}
\sum_{\substack{S\subseteq\{1,\ldots,n\}\\ |S|=r}}
\prod_{i\in S}(t-x_i).
\end{equation}
Equivalently,
\begin{equation}\label{eq:terminal-coefficients}
J_{n,r}(t)
=t^r+\sum_{j=2}^r(-1)^j
\frac{\binom rj}{\binom nj}e_jt^{r-j}
=t^r+\sum_{j=2}^r(-1)^j
\frac{(r)_{\downarrow j}}{(n)_{\downarrow j}}e_jt^{r-j}.
\end{equation}
Consequently $J_{n,r}$ is a universal polynomial in
$p_2,p_3,\ldots,p_r$, and
\[
\D_{n-r}^{(n)}=\left(\frac{n!}{r!}\right)^{2r-2}
\disc(J_{n,r}).
\]
Moreover, with $J_{n,1}(t)=t$ and $J_{n,0}(t)=1$, the terminal polynomials
form a finite Appell chain:
\begin{equation}\label{eq:appell}
J'_{n,r}(t)=rJ_{n,r-1}(t),\qquad 1\le r\le n.
\end{equation}
In particular,
\begin{equation}\label{eq:terminal-resultant}
\disc(J_{n,r})
=(-1)^{r(r-1)/2}r^r\Res(J_{n,r},J_{n,r-1}).
\end{equation}
If the roots $x_i$ are real, then all $J_{n,r}$ are real-rooted, and the
zeros of consecutive members interlace.
\end{theorem}

\begin{proof}
Differentiate the product $P(t)=\prod_{i=1}^n(t-x_i)$ exactly $n-r$
times.  For each $r$-subset $S$, the product
$\prod_{i\in S}(t-x_i)$ occurs $(n-r)!$ times, whence
\[
P^{(n-r)}(t)=(n-r)!\sum_{|S|=r}\prod_{i\in S}(t-x_i).
\]
Multiplication by $r!/n!=\bigl(\binom nr(n-r)!\bigr)^{-1}$ gives
\eqref{eq:subset-average}.  The coefficient of $(-1)^je_jt^{r-j}$ in
the average is the probability that a fixed $j$-subset is contained in a
uniformly chosen $r$-subset, namely
$\binom{n-j}{r-j}/\binom nr=\binom rj/\binom nj$; this proves
\eqref{eq:terminal-coefficients}.  Since $e_1=p_1=0$, Newton's identities
express every $e_j$ as a polynomial in $p_2,\ldots,p_j$.

The discriminant scaling formula gives the displayed normalization of
$\D_{n-r}^{(n)}$.  Differentiating the definition of $J_{n,r}$ yields
\[
J'_{n,r}=\frac{r!}{n!}P^{(n-r+1)}=rJ_{n,r-1},
\]
which is \eqref{eq:appell}.  Since $J_{n,r}$ is monic,
\[
\disc(J_{n,r})=(-1)^{r(r-1)/2}\Res(J_{n,r},J'_{n,r}),
\]
and \eqref{eq:terminal-resultant} follows from the homogeneity of the
resultant in its second argument.  Finally, repeated Rolle interlacing for
successive derivatives of a real-rooted polynomial gives real-rootedness and
interlacing of the entire chain.
\end{proof}


\end{document}
